\documentclass[11pt]{article}
\usepackage[T1]{fontenc}
\usepackage[utf8]{inputenc}
\usepackage{lmodern}
\usepackage[margin=1in]{geometry}
\usepackage{amsmath,amssymb,amsthm,mathtools,microtype,booktabs}
\usepackage[hidelinks]{hyperref}
\newtheorem{theorem}{Theorem}[section]
\newtheorem{lemma}[theorem]{Lemma}
\newtheorem{proposition}[theorem]{Proposition}
\newtheorem{corollary}[theorem]{Corollary}
\newcommand{\Pow}{\mathcal P_+}
\newcommand{\ow}[1]{\mathbin{\star_{#1}}}
\newcommand{\Sym}{\operatorname{Sym}}
\newcommand{\FP}{\mathcal{FP}^{+}}
\title{Owned Choice for Arbitrarily Many Roles:\\Ambiguity Actions and Equational Complexity}
\author{KFOL research working note}
\date{28 October 2026}
\begin{document}
\maketitle
\begin{abstract}
The rational owned-choice algebra has an exponential carrier but optimal
singleton-and-pair tests for its equations. We identify the ambiguity-action
quotient of its unary term monoid for every number of roles. It is a
quotient of an affine extension of the known factorpower of the symmetric
group. Its canonical rows satisfy an exact perfect-matching criterion;
recognition, synthesis, and composition are polynomial in the number of
roles. The quotient has $34$ elements for three roles and $1979$ for four.
The partial-permutation description is precisely the three-role case.
Equational validity is coNP-complete already for each fixed number of
roles at least three, even without role switches. The units of the full
unary term monoid are exactly the role permutations for every number of
roles. These statements characterize a quotient and a decision problem;
the complete higher-role unary clone and a finite equational basis remain
open. The Boolean embedding used for hardness is kernel-checked in Lean;
the general matching and complexity theorems have written proofs and
independent finite checks.
\end{abstract}

\section{A small algebra and its questions}

Fix a finite role set $I$ of size $K\geq3$. Its values are nonempty
winner sets. In finite public games with one winner at each leaf, each
role values its own win at one and every other outcome at zero. The
possible winners over full pure subgame-perfect equilibria compose by
\[
 S\ow r T=\begin{cases}\{r\},&S=\{r\}\text{ or }T=\{r\},\\
 S\cup T,&\text{otherwise}.\end{cases}
\]
The intended signature also has singleton constants and internal role
switches $S\mapsto\pi S$ for every $\pi\in\Sym(I)$. These switches
rename owners and winners inside a component. Terms have the grammar
\[
 t::=\langle i\rangle\mid x_j\mid t\oplus_r t\mid\pi(t).
\]
The full-strategy interpretation and exact equilibrium recurrence are
proved in \texttt{Basic.lean}. Hidden choices and sequential equilibrium
are outside this public fragment.

Three distinct owners define neutral union:
\[
 J(x,y)=(x\ow r y)\ow t(x\ow s y)=x\cup y.
\]
Consequently every nonempty constant, including $I$, is ground definable.
We use these derived operations freely. Repeated variable occurrences
denote independently solved copies, as required by the term algebra.

The preceding note \texttt{symmetric.tex} proves two uniform facts used
here. On each fixed singleton-or-ambiguous status cell, every term is a
constant union permutation images of its ambiguous inputs. Every failed
equation therefore has a countervaluation in which each input has at
most two winners. Every singleton and every pair is necessary in a fixed
test domain. Those mathematical proofs, including the separating terms,
are not yet complete Lean developments.

The three-role unary classification in that note has a quotient by
partial permutations. For more roles, a coordinate can depend on two or
more input membership bits: $x\cup(1\ 2)(x)$ at $K=4$ is an elementary
example. The right general object retains these disjunctive source sets,
and also the constant coordinates.

\section{The known relation monoid and its affine extension}

Regard a relation $R\subseteq I\times I$ as sending source roles to
output roles, and write $R[S]$ for its direct image. Let $\FP(\Sym(I))$
be the monoid of nonempty unions of permutation graphs, with relational
composition. This is the \emph{factorpower} of the natural symmetric-group
action, introduced by Ganyushkin and Mazorchuk and surveyed by
Mazorchuk\cite{GanyushkinMazorchuk1995,Mazorchuk2009}. In the usual
definition, nonempty subsets $P,Q\subseteq\Sym(I)$ are identified when
$\{\pi(i):\pi\in P\}=\{\pi(i):\pi\in Q\}$ for every $i$.
Thus the relation monoid appearing here is exactly a known object.

Form its affine extension
\[
 \mathcal E_I=\mathcal P(I)\rtimes\FP(\Sym(I)),\qquad
 (C,R)(D,Q)=(C\cup R[D],R\circ Q).
\]
The action on $\mathcal P(I)$ is by direct image; it preserves empty set
and unions, so this multiplication is associative. Its identity is
$(\varnothing,\mathrm{id})$. The affine action is
\begin{equation}\label{eq:affine}
 F_{C,R}(S)=C\cup R[S].
\end{equation}
Every such map sends ambiguous sets to ambiguous sets because $R$
contains a permutation. Restrict these maps to
$\mathcal A_I=\{S\subseteq I:|S|\geq2\}$ and identify equal restrictions.
Call the resulting transformation monoid $\mathcal Q_I$.

\begin{theorem}[Exact ambiguity quotient]\label{thm:quotient}
Let $\mathcal M_I$ be the monoid of unary term functions fixing $I$.
Restriction to $\mathcal A_I$ is a surjective monoid homomorphism
$\mathcal M_I\twoheadrightarrow\mathcal Q_I$.
\end{theorem}
\begin{proof}
A unary term is either globally constant or has live provenance
$C\cup\bigcup_{\pi\in P}\pi S$ on ambiguous inputs, with $P$ nonempty.
Here is a simultaneous induction establishing that dichotomy. The
variable is live and a winner is constant. Switching transports constants
and permutations. At an owned node, a live child is always ambiguous on
ambiguous inputs. Thus collapse there occurs exactly when a globally
constant child is the owner singleton, making the whole node globally
constant. Otherwise union the live provenances and the constant children.
A node with only constant children is again globally constant.

Every live expression fixes $I$. Hence a top-preserving unary term is
live or is the constant $I$, and has a restriction of the displayed form.
Conversely, for any $(C,R)$ choose a nonempty finite set of permutations
whose union is $R$. The term joining their images of $x$ and the
singleton constants in $C$ realizes \eqref{eq:affine} on all inputs and
fixes $I$. Restriction respects composition since all these maps preserve
$\mathcal A_I$.
\end{proof}

The constant part cannot be suppressed. For example,
$S\mapsto S\cup\{a\}$ on ambiguous inputs cannot equal $R[S]$ for a
union of permutations $R$. Every other output coordinate $b\ne a$ has
the singleton source $b$. As shown below, that source is uniquely
determined by the restricted action. Thus every constituent permutation
must fix all $b\ne a$, hence also $a$. Their union is the identity and
cannot make $a$ constantly present. The relevant object is an affine
extension and a quotient of a known monoid, not just the factorpower.

\section{Canonical rows, matching, and synthesis}

For an output role $k$ let
$R^{-1}(k)=\{i:(i,k)\in R\}$. Membership of $k$ in
\eqref{eq:affine} is either constant or the predicate
$S\cap R^{-1}(k)\ne\varnothing$. Every such source set is nonempty.
On ambiguous inputs this predicate is constantly true exactly when its
source set has size at least $K-1$.

\begin{lemma}[Unique row representation]\label{lem:rows}
Every action in $\mathcal Q_I$ has a unique row for each output role,
of one of the following forms:
\[
 *\quad\text{(always present)},\qquad
 A_k\subseteq I,\quad 1\leq|A_k|\leq K-2
 \quad\text{(present iff }S\cap A_k\ne\varnothing\text{)}.
\]
\end{lemma}
\begin{proof}
The preceding threshold gives existence. A nonconstant row has a pair
disjoint from its source set, distinguishing it from $*$. If two small
source sets $A,B$ differ, choose $i\in A\setminus B$, swapping their
names if necessary. Since $|B|\leq K-2$, choose a different
$j\notin B$. The pair $\{i,j\}$ meets $A$ and misses $B$, distinguishing
the rows.
\end{proof}

Not every syntactically possible collection of rows is realizable. Build
a bipartite graph $H$ with source and output copies of $I$. A nonconstant
output row $k$ allows exactly the edges from $A_k$; a constant row allows
edges from every source. Edges into nonconstant rows are
\emph{prescribed}; edges into constant rows are only allowed.

\begin{theorem}[Matching criterion]\label{thm:matching}
A collection of canonical rows belongs to $\mathcal Q_I$ exactly when
every prescribed edge belongs to a perfect matching of $H$. If all rows
are constant, it always belongs.
\end{theorem}
\begin{proof}
For necessity, choose a representative $(C,R)$. A nonconstant row
cannot belong to $C$, and its preimage under $R$ is exactly its canonical
source set by Lemma~\ref{lem:rows}. Every prescribed edge lies in some
permutation constituent of $R$. That permutation is a perfect matching
of $H$.

Conversely, choose a perfect matching for each prescribed edge. Their
union $R$ is a union of permutations. At a nonconstant row its preimage
is contained in the allowed source set and includes every prescribed
source, so it is exactly $A_k$. Set $C$ equal to the set of constant
output rows. Then $(C,R)$ realizes the given rows. If no edges are
prescribed, take $R$ to be the identity and $C=I$.
\end{proof}

For comparison, the analogous criterion with \emph{all} relation edges
prescribed is the standard factorpower realizability criterion
(Mazorchuk\cite[Section~2.1]{Mazorchuk2009}). Here the constant rows leave
their matching edges free, and the ambiguity threshold removes larger
nonconstant source sets. An ordinary Hall existence test alone is
insufficient. At $K=4$, rows
\[
 (\{1\},\{1,2\},*,*)
\]
allow a perfect matching, but the prescribed edge from source $1$ to
output $2$ belongs to none: output $1$ already requires source $1$.
The collection is therefore not realizable.

\begin{corollary}[Polynomial recognition and synthesis]
Given $K$ canonical rows, membership in $\mathcal Q_I$ is decidable in
polynomial time. Every member has a realization by at most $K^2$
permutation images, together with its constant roles, and by an ordinary
term of size polynomial in $K$.
\end{corollary}
\begin{proof}
For each of the at most $K(K-2)$ prescribed edges, force that edge and
find a perfect matching in the remaining graph. Standard augmenting
paths suffice. Use those witness permutations in the synthesis from
Theorem~\ref{thm:quotient}. Join the pieces in a balanced tree. Expanding
the fixed derived-union term duplicates each argument twice; balanced
joining of $m$ pieces has size $O(m^2)$ rather than the exponential
expansion of a skew join. Singleton constant pieces and explicit
permutation descriptions have polynomial size too.
\end{proof}

\begin{theorem}[Canonical composition]\label{thm:composition}
For outer rows $A$ and inner rows $B$, the composite row at $k$ is:
\begin{enumerate}
\item $*$ if $A_k=*$ or some $i\in A_k$ has $B_i=*$;
\item otherwise $\bigcup_{i\in A_k}B_i$, changed to $*$ if that union
has at least $K-1$ elements.
\end{enumerate}
This computes composition in time polynomial in $K$ on descriptions of
$O(K^2)$ bits.
\end{theorem}
\begin{proof}
Substitute the inner membership disjunctions into the outer membership
disjunction. A constant true disjunct makes the result true; otherwise
the sources union. Apply the ambiguity threshold. Closure follows also
from the affine multiplication in $\mathcal E_I$.
\end{proof}

For three roles every nonconstant row is a singleton. The matching
criterion says that their sources must be distinct. Thus these rows
are exactly a partial bijection, with $*$ on the missing output roles,
recovering the $34$-element symmetric inverse monoid. Four roles admit
two-source rows, and the matching criterion is a genuine further
constraint. Exhausting the criterion and, independently, all affine
permutation-union representatives gives
\[
\begin{array}{c|rr}
K&|\FP(\Sym(I))|&|\mathcal Q_I|\\\hline
3&49&34\\
4&7443&1979
\end{array}
\]
These are finite computations, with complete enumerations reproduced by
\texttt{check\_general.py}. They count the ambiguity quotient, not the
full unary term monoid. Its singleton fibers for general $K$ are still
unclassified.

\section{Reversible contexts at every number of roles}

\begin{lemma}
Every term function is monotone for inclusion on nonempty winner sets.
\end{lemma}
\begin{proof}
Constants, variables, and permutation images are monotone. Suppose
$S\subseteq S'$ and $T\subseteq T'$, all nonempty. If an upper operand
equals $\{r\}$, its corresponding lower operand is also $\{r\}$, and
both owned outputs equal $\{r\}$. Otherwise the upper output is
$S'\cup T'$. The lower output is either $S\cup T$, which is contained
in it, or $\{r\}$ supplied by one lower operand, which is contained in
the corresponding upper operand. Thus owned choice is monotone.
\end{proof}

\begin{theorem}[All unary units]
For every $K\geq3$, the group of units of the full unary term monoid is
exactly $\Sym(I)$.
\end{theorem}
\begin{proof}
A unit and its term inverse are monotone, so they form an order
automorphism of the poset $\Pow(I)$. Such an automorphism permutes its
minimal elements, the singletons. Every set is determined by the
singletons below it, so that permutation determines the whole map.
Conversely every role permutation and its inverse are terms.
\end{proof}

This generalizes the three-role unit calculation without requiring a
higher-role classification of singleton fibers. It is a direct use of
the familiar order-automorphism argument for the Boolean subset poset.

\section{Equational validity is coNP-complete}

The decision problem takes two ordinary tree terms, with any number of
variables, and asks whether they agree at every nonempty-set valuation.
Fixing $K$ fixes the algebra and its signature. In the uniform version,
role names have binary encodings and permutations are explicit lists of
their images. A succinct encoding of terms by shared recursive
descriptions is not assumed.

\begin{theorem}[Equational complexity]\label{thm:complexity}
For every fixed $K\geq3$, equational validity is coNP-complete, already
in the signature without role switches. The uniform problem with $K$
as part of the input is coNP-complete as well.
\end{theorem}
\begin{proof}
For fixed $K$, a failed valuation is a polynomial certificate since the
algebra is finite and term evaluation is linear up to fixed constants.
For variable $K$, use the pair countervaluation theorem. A certificate
lists at most two role labels per input variable. Evaluate terms using
sparse sets: an owned node either collapses or unions its inputs, while
a permutation preserves cardinality. Each intermediate set has size at
most twice the number of variable leaves below it plus the number of
winner leaves. This is polynomial in tree-term size, independently of
the exponential size of the carrier and even when $K$ is binary encoded.
Evaluation and comparison are therefore polynomial. Validity is in coNP.

For hardness fix distinct roles $a,b$ and encode Boolean truth by
$\mathbf1=\{a\}$ and falsehood by $\mathbf0=I$. The gate
\[
 q(x)=x\ow a I
   =\begin{cases}\mathbf1,&x=\{a\},\\\mathbf0,&\text{otherwise}
     \end{cases}
\]
is a surjection onto these two values. On them,
\[
 u\ow a v=u\lor v,\qquad u\ow b v=u\land v.
\]
The first owner accepts either true operand. For the second owner no
operand is its singleton, so it takes union; union absorbs the false
value $I$. Consequently any monotone Boolean formula equation can be
translated, replacing every variable by $q(x)$, disjunction by $\star_a$,
and conjunction by $\star_b$. It is valid in the owned-choice algebra
exactly when the Boolean equation is valid. The term for $I$ has fixed
size because $K$ is fixed. This is a polynomial translation using no
role switches.

The standard dual-rail reduction gives hardness directly. Given a CNF
formula $\varphi(z_1,\ldots,z_n)$, replace positive occurrences
of $z_i$ by $p_i$ and negative occurrences by $n_i$, producing a monotone
formula $\Phi(p,n)$. Put
\[
 A=\left(\bigwedge_i(p_i\lor n_i)\right)\land\Phi(p,n),
 \qquad B=\bigvee_i(p_i\land n_i).
\]
The equation $A\lor B=B$ fails exactly when $A=1$ and $B=0$.
Those conditions choose exactly one rail for every variable and satisfy
$\varphi$. Conversely any satisfying assignment provides such rails.
Thus $\varphi$ is unsatisfiable exactly when the equation is valid.
CNF unsatisfiability is coNP-complete by the standard SAT theorem
\cite{Cook1971}. Formula sizes and their owned-choice translations are
polynomial. Hardness at the fixed value $K=3$ also gives uniform hardness.
\end{proof}

The two-value encoding does not replace the winner-set algebra by a
Boolean algebra. It is a selected-operation fragment with a definable
gate. In particular it is not a subalgebra of the full signature with
all named winner constants and all internal role permutations. The
complexity method is standard; its consequence here closes the lower
bound left open by pair testing.

\section{Evidence, attribution, and the next mathematical gap}

\texttt{BooleanCore.lean} proves the gate, the two Boolean operations,
and the embedding of arbitrary monotone Boolean formula equations. It
uses ordinary syntax from the preceding modules and proves both
directions of validity preservation for all nonempty input valuations.
The SAT reduction and complexity-class statements remain mathematical
proofs above. Lean 4.34.1 checks the module with only \texttt{Std}; the
axiom audit reports only \texttt{propext}, \texttt{Classical.choice}, and
\texttt{Quot.sound}, with no admitted statements or new axioms.

\texttt{ambiguity.py} implements the matching criterion, witness
permutations, balanced syntax synthesis, and canonical composition.
\texttt{check\_general.py} compares two independent complete
enumerations at $K=3,4$: admissible canonical rows versus affine images
of the factorpower, itself enumerated by union closure of all
permutations. It checks semantic uniqueness, every composition at $K=3$,
sampled compositions at $K=4$, and actual term synthesis. It also checks
the Boolean gates through $K=6$ and the dual-rail equivalence on all
$4369$ CNFs with up to three clauses on two variables, including empty
clauses. These support the written proofs, rather than replacing them.

The factorpower, its edge-extension criterion, semidirect-product
construction, matching algorithms, Boolean hardness reduction, and
subset-poset automorphism argument are established mathematics reused
explicitly. The candidate contribution is their exact relationship to
rational owned choice: the optimal pair theorem from the preceding note,
the affine ambiguity quotient with its threshold and canonical matching
criterion, and their consequences for contexts and equational validity.
Priority of this particular combination is not established by identifying
the factorpower alone.

There is still no finite equational basis or proof calculus. The
all-arity pair theorem controls input tests, but not which tables on
those tests are realizable. The ambiguity quotient controls one part of
unary composition, but forgets singleton rows. Two focused next questions
are therefore the exact singleton fibers at general $K$, and a uniform
presentation that derives the status split equationally. The residual
bound needed by the conditional Willard route in \texttt{symmetric.tex}
also remains open. No dependent type-theoretic claim follows yet.

\bibliographystyle{plain}
\bibliography{references}
\end{document}
